\documentclass{article}
\usepackage[margin=1in]{geometry} % set suitable margins
\usepackage{parskip}
\usepackage{graphicx}

\usepackage{amsmath}
\numberwithin{equation}{section}

\title{EMCH 354 HW 5 Solutions}
\author{Jacob Vaught}
\date{}

\begin{document}
\maketitle

\section{Problem \#1}
\subsection*{Given Data}
\begin{itemize}
    \item Diameter of the copper sphere, $d = 2 \, \text{cm} = 0.02 \, \text{m}$
    \item Maximum rate of temperature change, $\frac{dT}{dt} = 19.8 \, \text{K/s}$
    \item Temperature of the sphere when maximum rate of temperature change occurs, $T = 92.5 \, \text{K}$
    \item Saturation temperature of nitrogen at $1 \, \text{atm}$ pressure, $T_{\infty} = 77.4 \, \text{K}$
    \item Density of copper, $\rho = 8930 \, \text{kg/m}^3$
    \item Specific heat capacity of copper, $c = 235 \, \text{J/(kg} \cdot \text{K)}$
    \item Thermal conductivity of copper at $92.5 \, \text{K}$, $k = 450 \, \text{W/(m} \cdot \text{K)}$
\end{itemize}

\subsection*{Objective}
\begin{enumerate}
    \item Determine the heat transfer coefficient using the lumped thermal capacity model.
    \item Check the Biot number to validate the use of the lumped model.
\end{enumerate}

\subsection*{Solution}

\subsubsection*{Part 1: Determining the Heat Transfer Coefficient}
The lumped thermal capacity model is based on the premise that the temperature within the solid varies negligibly, allowing the assumption of uniform temperature. The rate of heat transfer, $Q$, can be expressed in two ways:
\begin{equation}
    Q = mc\frac{dT}{dt} = hA(T - T_{\infty})
\end{equation}
where $m$ is the mass of the copper sphere, $c$ is the specific heat capacity, $\frac{dT}{dt}$ is the rate of temperature change, $h$ is the heat transfer coefficient, $A$ is the surface area, $T$ is the temperature of the sphere, and $T_{\infty}$ is the ambient temperature.
\newline\newline
The mass of the sphere can be calculated using its volume and the density of copper:
\begin{align}
    V &= \frac{4}{3}\pi (0.01 \, \text{m})^3 \\
    m &= \rho V = 8930 \, \text{kg/m}^3 \times \frac{4}{3}\pi (0.01 \, \text{m})^3
\end{align}
\noindent
The surface area of the sphere is given by:
\begin{equation}
    A = 4\pi (0.01 \, \text{m})^2
\end{equation}
\noindent
Substituting these expressions into the equation for $Q$, and solving for $h$, we find:
\begin{equation}
    h = \frac{mc\frac{dT}{dt}}{A(T - T_{\infty})} = \frac{8930 \, \text{kg/m}^3 \times \frac{4}{3}\pi (0.01 \, \text{m})^3 \times 235 \, \text{J/(kg} \cdot \text{K)} \times 19.8 \, \text{K/s}}{4\pi (0.01 \, \text{m})^2 \times (92.5 \, \text{K} - 77.4 \, \text{K})}
\end{equation}
\noindent
Using the given values, the heat transfer coefficient is calculated to be approximately $9172.47 \, \text{W/(m}^2\text{K)}$.

\subsubsection*{Part 2: Validity Check using the Biot Number}
The Biot number, $Bi$, is a dimensionless number that compares internal resistance to heat conduction within the body to the resistance to heat transfer across the boundary. It is defined as:
\begin{equation}
    Bi = \frac{hL_c}{k}
\end{equation}
where $L_c = \frac{V}{A}$ is the characteristic length of the sphere, and $k$ is the thermal conductivity of the material. For our sphere, $L_c$ is calculated as the volume divided by the surface area.
\newline\newline
The Biot number, using the calculated heat transfer coefficient, is:
\begin{equation}
    Bi = \frac{9172.47 \, \text{W/(m}^2\text{K)} \times \frac{\frac{4}{3}\pi (0.01 \, \text{m})^3}{4\pi (0.01 \, \text{m})^2}}{450 \, \text{W/(m} \cdot \text{K)}}
\end{equation}
\noindent
This calculation yields a Biot number of approximately $0.0679$. Since $Bi << 0.1$, the lumped thermal capacity model is valid, indicating that the temperature gradient within the sphere is negligible.

\newpage
\section{Problem \#2}
\subsection*{Given Data}
Electronic components are mounted on a finned aluminum base plate with effective thermal conduction paths, exposed to a cooling air stream from a fan. The system's characteristics are as follows:
\begin{itemize}
    \item The sum of the products of mass and specific heat of the base plate and components, $mc = 5000 \, \text{J/K}$
    \item The product of heat transfer coefficient and surface area, $hA = 10 \, \text{W/K}$
    \item Initial temperature of the plate and cooling air, $T_{\infty} = 295 \, \text{K}$
    \item Power applied to the system, $Q = 300 \, \text{W}$
\end{itemize}

\subsection*{Objective}
\begin{enumerate}
    \item Determine the plate temperature after 10 minutes.
    \item Find the steady-state temperature of the components.
\end{enumerate}

\subsection*{Solution}

\subsubsection*{Part A}
The temperature of the plate at any time, $T(t)$, can be described by the transient solution:
\begin{equation}
    T(t) = T_{\infty} + (T_{ss} - T_{\infty})(1 - e^{-\frac{t}{\tau}})
\end{equation}
where $\tau$ is the time constant of the system, defined as:
\begin{equation}
    \tau = \frac{mc}{hA} = \frac{5000 \, \text{J/K}}{10 \, \text{W/K}} = 500 \, \text{s}
\end{equation}
For $t = 10$ minutes ($600$ seconds), the temperature is:
\begin{equation}
    T(600) = 295 \, \text{K} + (325 \, \text{K} - 295 \, \text{K})(1 - e^{-\frac{600}{500}}) \approx 315.96 \, \text{K}
\end{equation}

\subsubsection*{Part B}
In the steady state, the heat generated by the electronic components is equal to the heat dissipated to the surrounding air. The steady-state temperature, $T_{ss}$, can be found using the balance equation:
\begin{equation}
    Q = hA(T_{ss} - T_{\infty})
\end{equation}
Solving for $T_{ss}$ gives:
\begin{equation}
    T_{ss} = \frac{Q}{hA} + T_{\infty} = \frac{300 \, \text{W}}{10 \, \text{W/K}} + 295 \, \text{K} = 325 \, \text{K}
\end{equation}

\newpage
\section{Problem \#3}
Consider a 10 m length of 2 cm internal diameter (I.D.) tube. The tube wall is maintained at a temperature of 320 K, and water enters the tube at 300 K, with an atmospheric pressure of 1 atm and a flow velocity of 3 m/s. Determine the average heat transfer coefficient and the pressure gradient inside the tube under these conditions.

\subsection*{Solution}

\textbf{Bulk Mean Temperature Calculation:}
The bulk mean temperature \( T_b \) is calculated as the average of the entering water temperature and the tube wall temperature. This temperature is representative of the fluid's core temperature.
\begin{equation}
T_b = \frac{T_{\text{inlet}} + T_{\text{wall}}}{2} = \frac{320 + 300}{2} = 310\, \text{K}
\end{equation}

\textbf{Property Evaluation at \( T_b = 310\, \text{K} \):}
The thermophysical properties of water are evaluated at the bulk mean temperature, \( T_b \), to ensure accuracy in the subsequent calculations.
\begin{equation}
\left\{
\begin{array}{l}
L = 10\, \text{m} \quad \text{(length of the tube)}, \\
D = 2\, \text{cm} \quad \text{(internal diameter of the tube)}, \\
V_0 = 3\, \text{m/s} \quad \text{(velocity of water entering the tube)}, \\
\text{Pr} = 4.6 \quad \text{(Prandtl number, a dimensionless number relating the momentum diffusivity to thermal diffusivity)}, \\
k = 0.628\, \text{W/mK} \quad \text{(thermal conductivity of the fluid)}, \\
\rho = 993\, \text{kg/m}^3 \quad \text{(density of the fluid)}, \\
\mu = 6.95 \times 10^{-4}\, \text{Pa}\cdot\text{s} \quad \text{(dynamic viscosity of the fluid)}, \\
c_p = 4174\, \text{J/(kgK)} \quad \text{(specific heat capacity at constant pressure)}
\end{array}
\right.
\end{equation}

\textbf{Reynolds Number Calculation:}
The Reynolds number \( Re_D \) is a dimensionless number that helps predict flow patterns in different fluid flow situations. In this case, it's used to determine whether the flow is laminar or turbulent.
\begin{equation}
Re_D = \frac{\rho V_0 D}{\mu} = \frac{993 \times 3 \times 0.02}{6.95 \times 10^{-4}} = 85726.6 \quad (\text{indicating turbulent flow since } Re_D > 2300)
\end{equation}

\textbf{Heat Transfer and Flow Characteristics:}
Using empirical correlations, we can estimate the Nusselt number (Nu) for heat transfer and the friction factor (f) for fluid flow in turbulent conditions.
\begin{equation}
f_b = \left(0.79 \ln(\text{Re}_D) - 1.64\right)^{-2} = \left(0.79 \ln(85726.6) - 1.64\right)^{-2} = 0.0186
\end{equation}

\begin{equation}
Nu_b = \frac{\frac{f}{8}(\text{Re}_D - 1000)\text{Pr}}{1 + 12.7 \left(\frac{f}{8}\right)^{0.5}(\text{Pr}^{2/3} - 1)} = \frac{\frac{0.0186}{8}(85726 - 1000)4.6}{1 + 12.7 \left(\frac{0.0186}{8}\right)^{0.5}(4.6^{2/3} - 1)} = 435.25
\end{equation}

\textbf{Temperature Correction Factors:}
Due to the variable properties of the fluid at different temperatures, we apply corrections to account for the changes in viscosity at the surface and bulk mean temperatures.
\begin{equation}
\frac{\mu_s}{\mu_b} = \text{(viscosity correction factor)} = \frac{\mu_{\text{at surface temperature}}}{\mu_{\text{at bulk mean temperature}}} = \frac{5.84 \times 10^{-4}}{6.95 \times 10^{-4}} = 0.84
\end{equation}

\textbf{Corrected Heat Transfer Coefficient and Friction Factor:}
The final calculations adjust the heat transfer coefficient (h) and pressure gradient (\(\Delta P/L\)) to account for temperature effects on viscosity.
\begin{equation}
f = f_b \left(\frac{\mu_s}{\mu_b}\right)^{0.25} = 0.0186 \times 0.84^{0.25} = 0.0178
\end{equation}
\begin{equation}
Nu = Nu_b \left(\frac{\mu_s}{\mu_b}\right)^{-0.11} = 435.25 \times 0.84^{-0.11} = 443.7
\end{equation}

\textbf{Final Results:}
Using the corrected Nusselt number and thermal conductivity, we calculate the heat transfer coefficient. The friction factor is used along with the density and velocity to find the pressure gradient.
\begin{equation}
h = \frac{Nu \times k}{D} = \frac{443.7 \times 0.628}{0.02} = 13932\, \frac{\text{W}}{\text{m}^2 \text{K}}
\end{equation}
\begin{equation}
\frac{\Delta P}{L} = \frac{1}{2} f \rho V_0^2 = \frac{1}{2} \times 0.0178 \times 993 \times 3^2 = 3977.2\, \text{Pa/m}
\end{equation}


\newpage
\section{Problem \#4} 
A very thin stainless steel shim stock is bonded to a 10 cm-square Teflon plate to serve as a heater element. Water at 290 K flows as a laminar boundary along the plate. The voltage drop and current flow in the heater are measured to obtain the power dissipated; also, from the known dependence of electrical resistance on temperature, the average plate temperature can be deduced. In a particular test, the power is 1025 W, and the average plate temperature is 311.5 K. Estimate the water velocity, and the temperature distribution along the plate.

% Dimensions and thermal properties
\noindent
First, we start by defining the dimensions and thermal properties of the system:

\begin{equation}
L = 10 \text{ cm} = 0.1 \text{ m} \quad \text{(length of the plate)}
\end{equation}
\begin{equation}
T_{\infty} = 290 \text{ K} \quad \text{(temperature of the water flowing over the plate)}
\end{equation}
\begin{equation}
T_{s} = 311.5 \text{ K} \quad \text{(average surface temperature of the plate)}
\end{equation}
\begin{equation}
\text{Average temperature, } T_{r} = \frac{290 + 311.5}{2} = 300.8 \text{ K} \quad \text{(used for some thermal calculations)}
\end{equation}
\begin{equation}
Pr = 5.78 \quad \text{(Prandtl number, a dimensionless number relevant in fluid flow)}
\end{equation}
\begin{equation}
k = 0.6125 \text{ W/m.K} \quad \text{(thermal conductivity of the fluid)}
\end{equation}
\begin{equation}
\nu = 8.54 \times 10^{-7} \text{ m}^2/\text{s} \quad \text{(kinematic viscosity of the fluid)}
\end{equation}

% Assumptions
\noindent
We assume a laminar flow without significant conduction in the plate's direction or through its backside, as it is insulated. This simplifies our calculations to focus on heat transfer directly from the fluid to the plate's surface:

% Heat transfer calculations
\begin{equation}
Nu_{L} = 0.68 Re_{L}^{\frac{1}{2}} Pr^{\frac{1}{3}} = \frac{h_{c} \times L}{k} \quad \text{(defining Nusselt number using the heat transfer coefficient, \( h_c \))}
\end{equation}
\begin{equation}
Nu_{L} = \frac{h_{c} \times L}{k} = \frac{\frac{q_{s}}{T_{s} - T_{\infty}} \times L}{k} = \frac{\frac{Q}{A} \times L}{T_{s} - T_{\infty}} \times \frac{L}{k} = \frac{\left( \frac{1025}{(0.1 \times 0.1)} \right) \times 0.1}{(311.5 - 290)} \times \frac{0.1}{0.615} = 778.4
\end{equation}

% Flow characteristics
\begin{equation}
Nu_{L} = 0.68 Re_{L}^{\frac{1}{2}} Pr^{\frac{1}{3}} = 778.4 \quad \text{(calculation from earlier, used to find Reynolds number)}
\end{equation}
\begin{equation}
Re_{L} = \left( \frac{778.4}{0.68 \times Pr^{\frac{1}{3}}} \right)^{2} = 406807 = 4.06 \times 10^{5} \quad \text{(Reynolds number indicating laminar flow)}
\end{equation}
\begin{equation}
Re_{L} = \frac{V \times L}{\nu} = 4.06 \times 10^{5} \times 8.54 \times 10^{-7}
\end{equation}
\begin{equation}
V = \frac{3.4742 \text{ m/s}}{0.1} = 3.4742 \text{ m/s} \quad \text{(water velocity calculated from Reynolds number)}
\end{equation}

% Temperature distribution
\noindent
Finally, we calculate the temperature distribution along the plate using the heat flux and previously determined flow characteristics:

\begin{equation}
Nu_{x} = \frac{h_{x} \times x}{k} = \frac{q'' \times x}{\Delta T_{x}} = \frac{q'' \times x}{k \Delta T_{x}} = \frac{q'' \times x}{k(T_{x} - T_{\infty})}
\end{equation}
\begin{equation}
\frac{1}{Re_{x}^{\frac{1}{2}}} = \frac{V \times x}{\nu}
\end{equation}
\begin{equation}
Nu_{x} = 0.453 Re_{x}^{\frac{1}{2}} Pr^{\frac{1}{3}} \Rightarrow T_{x} - T_{\infty} = \frac{q'' \times x^{\frac{1}{2}}}{0.453 k \times V^{\frac{1}{2}} \times Pr^{-\frac{1}{3}}}
\end{equation}
\begin{equation}
T_{x} = 102.058 \, x^{\frac{1}{2}} + 290 \quad \text{(expression for temperature distribution along the plate)}
\end{equation}
\noindent
Below is the graphical representation of the findings:
\begin{center}
    \includegraphics[width=0.7\textwidth]{image.png}
\end{center}
\newpage
\section{Problem \#5}

Consider the problem of SAE 50 oil flowing through a long tube. In many industrial applications, understanding the temperature profile of a fluid along a pipe is crucial for process control. This problem will address how to determine the exit temperature of the oil given specific conditions.

\textbf{Problem Statement:} SAE 50 oil flows in a 60 m-long, 25 mm-ID tube at 0.5 kg/s. The oil enters the tube at 300 K, and the walls of the tube are maintained at 370 K. The objective is to determine the exit bulk temperature of the oil.

\subsection*{Given Data}
The following parameters are given, which describe the physical setup and initial conditions of the problem:
\begin{align}
    \dot{m} &= 0.5 \, \text{kg/s} \quad &\text{(Mass flow rate of the oil)} \\
    D &= 25 \, \text{mm} = 0.025 \, \text{m} \quad &\text{(Internal diameter of the tube)} \\
    L &= 60 \, \text{m} \quad &\text{(Length of the tube)} \\
    T_{\text{in}} &= 300 \, \text{K} \quad &\text{(Inlet temperature of the oil)} \\
    T_{s} &= 370 \, \text{K} \quad &\text{(Temperature of the tube walls)}
\end{align}

\subsection*{Temperature Assumption}
To solve this problem, we first make an educated guess about the outlet temperature of the oil. We assume it to be 320 K and will verify this assumption at the end of our calculations:
\begin{equation}
    \frac{T_{\text{in}} + T_{\text{out}}}{2} = \frac{300 + 320}{2} = 310 \, \text{K}
\end{equation}

\subsection*{Fluid Properties at Reference Temperature}
The properties of the SAE 50 oil at the reference temperature, which is the average of the inlet and outlet temperatures, are given by:
\begin{align}
    Pr &= 3400 \quad &\text{(Prandtl number)} \\
    C_{p} &= 1950 \, \text{J/kg} \cdot \text{k} \quad &\text{(Specific heat capacity)} \\
    \rho &= 877 \, \text{kg/m}^3 \quad &\text{(Density)} \\
    k &= 0.1435 \, \text{W/m} \cdot \text{K} \quad &\text{(Thermal conductivity)} \\
    \mu &= 2500 \times 10^{-4} \, \text{kg/m} \cdot \text{s} \quad &\text{(Dynamic viscosity)}
\end{align}

\subsection*{Reynolds Number Calculation}
The Reynolds number is a dimensionless quantity used in fluid mechanics to characterize the flow regime of the fluid, indicating whether the flow is laminar or turbulent. It is calculated as follows:
\begin{equation}
    \text{Re}_{D} = \frac{\rho v D}{\mu} = \frac{\dot{m} D}{\pi D^2 \mu / 4} = \frac{4 \dot{m}}{\pi D \mu} = \frac{4 \times 0.5}{\pi \times 0.025 \times 2500 \times 10^{-4}} = 101.86 < 2300
\end{equation}

A Reynolds number below 2300 indicates laminar flow, and thus we confirm that the flow of oil in the tube is laminar.

\subsection*{Nusselt Number for Laminar Flow}
The Nusselt number, which is a dimensionless temperature gradient, is given for laminar flow inside a tube by the following correlation:
\begin{equation}
    Nu_{D} = 3.66 + \frac{0.065 \frac{D}{L} Re_{D} Pr}{1 + 0.04 \left(\frac{D}{L} Re_{D} Pr\right)^{\frac{1}{3}}} = 3.66 + \frac{0.065 \times \frac{0.025}{60} \times 101.86 \times 3400}{1 + 0.04 \left(\frac{0.025}{60} \times 101.86 \times 3400\right)^{\frac{1}{3}}} = 8.125
\end{equation}

\subsection*{Property Correction Factor}
Given the variation in fluid properties due to temperature differences, a correction factor for the Nusselt number must be applied. This is calculated using the viscosity at the wall and bulk temperatures:
\begin{equation}
    Nu_{D} = Nu_{D} \left(\frac{\mu_{s}}{\mu_{b}}\right)^{-0.11} = 8.125 \left(\frac{181}{2500}\right)^{-0.11} = 10.846
\end{equation}

\subsection*{Convective Heat Transfer Coefficient}
The convective heat transfer coefficient is a measure of the convective heat transfer between the fluid and the tube wall. It is derived from the Nusselt number:
\begin{equation}
    h_{c} = \frac{k}{D} Nu_{D} = \frac{0.1435}{0.025} \times 10.846 = 62.257 \, \text{W/(m}^2 \text{K)}
\end{equation}

\subsection*{Energy Balance and Exit Temperature}
Applying the energy balance around the tube, we can derive the exit bulk temperature of the oil:
\begin{equation}
    T_{s} - T_{\text{out}} = \frac{\rho L h_{c}}{\dot{m} c_{p}} (T_{s} - T_{\text{in}}) \quad \Rightarrow \quad T_{\text{out}} = T_{s} - (T_{s} - T_{\text{in}})e
\end{equation}

Upon simplification, the final outlet temperature is found to be:
\begin{equation}
    T_{\text{out}} = 318.190 \, \text{K}
\end{equation}

The calculated exit temperature is close to our initial guess, validating our assumptions. This also indicates that the properties did not change significantly and reevaluation is not necessary.

\newpage
\section{Problem \#6}
Calculate the average heat transfer coefficient for a hailstone with a diameter of 6 mm falling at its terminal velocity through air at a temperature of 287 K and a pressure of 1 atm.

\subsection*{Given Data:}
\begin{itemize}
    \item Diameter of the hailstone (\(d\)): 6 mm
    \item Temperature of the air (\(T\)): 287 K
    \item Pressure of the air (\(P\)): 1 atm
    \item Density of the hailstone (\(\rho_s\)): 917 kg/m\(^3\) (assumed to be constant)
    \item Density of air (\(\rho\)): 1.2 kg/m\(^3\) (standard density at sea level)
\end{itemize}

\subsection*{Terminal Velocity Calculation:}
To begin, we calculate the terminal velocity (\(V_t\)) of the hailstone. The terminal velocity is the constant speed that a freely falling object eventually reaches when the resistance of the medium through which it is falling prevents further acceleration. For a sphere, the terminal velocity can be calculated using the balance of gravitational and drag forces, given by the equation:

\begin{equation}
V_t = \sqrt{\frac{4 d \rho_s g}{3 C_D \rho}}
\end{equation}

Here, \(g\) is the acceleration due to gravity, \(C_D\) is the drag coefficient, and \(\rho\) is the density of the fluid through which the object is falling. The drag coefficient for a sphere is approximately 0.47 in turbulent flow, which is a common condition for falling hailstones. By substituting the known values, we find the terminal velocity:

\begin{equation}
V_t = \sqrt{\frac{4 \times 0.006 \times 917 \times 9.81}{3 \times 0.47 \times 1.2}}
\end{equation}
\begin{equation}
V_t = 11.29 \text{ m/s}
\end{equation}

\subsection*{Reynolds Number Calculation:}
The Reynolds number (\(Re\)) is a dimensionless quantity that helps predict flow patterns in different fluid flow situations. It is defined as the ratio of inertial forces to viscous forces and is used to determine whether a flow will be laminar or turbulent. The equation for the Reynolds number is:

\begin{equation}
Re = \frac{V_t \cdot d}{\nu}
\end{equation}

Here, \(\nu\) is the kinematic viscosity of the fluid. At the given temperature and pressure conditions, the kinematic viscosity of air (\(\nu\)) is approximately \(15 \times 10^{-6} \text{ m}^2/\text{s}\). Plugging in the values, we obtain:

\begin{equation}
Re = \frac{11.29 \cdot 0.006}{15 \times 10^{-6}} = 4516
\end{equation}

This \(Re\) indicates that the flow is likely to be turbulent, as it is greater than the critical Reynolds number for spheres, which is around 300 to 400.

\subsection*{Nusselt Number Calculation:}
The Nusselt number (\(Nu\)) is a dimensionless number that gives the ratio of convective to conductive heat transfer across a boundary. The Nusselt number for a sphere in fluid flow can be calculated by using empirical correlations based on experimental data, such as:

\begin{equation}
Nu = 2 + (0.4 \cdot Re^{0.5} + 0.06 \cdot Pr^{0.67}) \cdot Pr^{0.4}
\end{equation}

The Prandtl number (\(Pr\)) is a dimensionless number that characterizes the relative thickness of the momentum and thermal boundary layers. Given that \(Pr\) for air at this temperature is 0.713, we find \(Nu\) as:

\begin{equation}
Nu = 2 + (0.4 \cdot 4516^{0.5} + 0.06 \cdot 4516^{0.67}) \cdot 0.713^{0.4} = 40.2
\end{equation}

\subsection*{Heat Transfer Coefficient Calculation:}
The heat transfer coefficient (\(h\)) is a measure of the heat transfer rate per unit area per unit temperature difference. It is directly related to the Nusselt number through the relationship:

\begin{equation}
Nu = \frac{h \cdot d}{k}
\end{equation}

Solving for \(h\) and using the calculated \(Nu\) and the known thermal conductivity of air (\(k\)), we have:

\begin{equation}
h = \frac{Nu \cdot k}{d}
\end{equation}

\begin{equation}
h = \frac{40.2 \cdot 0.0257}{0.006} = 172.19 \text{ W/m}^2\text{K}
\end{equation}

This concludes our calculation. The average heat transfer coefficient for a 6 mm-diameter hailstone falling at its terminal velocity through air at 287 K and 1 atm is \( 172.19 \) W/m\(^2\)K.

\newpage
\section{Problem \#7}

Consider a counterflow, concentric tube heat exchanger that is designed to heat water from 20 to 80 \( ^{\circ}\)C. Hot oil is used as the heating agent, which enters the annulus at 160 \( ^{\circ}\)C and leaves at 140 \( ^{\circ}\)C. The inner tube, through which water flows, is thin-walled with a diameter of 20 mm. The overall heat transfer coefficient is given as 500 W/m\(^2\)K, and the design calls for a total heat transfer rate of 3000 W.

\subsection*{a. Determining the Length of the Heat Exchanger}

To find the required length of the heat exchanger, we first need to calculate the logarithmic mean temperature difference (LMTD). The LMTD is a correction factor that accounts for the temperature variation across the heat exchanger and is defined as:

\begin{equation}
\Delta T_{lm} = \frac{\Delta T_{o} - \Delta T_{i}}{\ln\left(\frac{\Delta T_{o}}{\Delta T_{i}}\right)}
\end{equation}

where:
\begin{itemize}
    \item \( \Delta T_{o} \) is the temperature difference between the hot and cold fluids at the outlet.
    \item \( \Delta T_{i} \) is the temperature difference at the inlet.
\end{itemize}

Given:
\begin{itemize}
    \item Inlet temperature of hot oil, \( T_{h,i} = 160 \, ^{\circ}C \).
    \item Outlet temperature of hot oil, \( T_{h,o} = 140 \, ^{\circ}C \).
    \item Outlet temperature of water, \( T_{c,o} = 80 \, ^{\circ}C \).
    \item Inlet temperature of water, \( T_{c,i} = 20 \, ^{\circ}C \).
\end{itemize}

We can now calculate the temperature differences:
\begin{equation}
\Delta T_{i} = T_{h,i} - T_{c,o} = 160 - 80 = 80 \, ^{\circ}C
\end{equation}
\begin{equation}
\Delta T_{o} = T_{h,o} - T_{c,i} = 140 - 20 = 120 \, ^{\circ}C
\end{equation}

Substituting into the LMTD equation, we get:
\begin{equation}
\Delta T_{lm} = \frac{120 - 80}{\ln(120/80)} = 98.65 \, ^{\circ}C
\end{equation}

Using the heat transfer rate equation \( q = U \times A \times \Delta T_{lm} \), and solving for the area \( A \), we can find the length \( L \) since the area is also a function of length (\( A = \pi D_i L \)):
\begin{equation}
L = \frac{q}{U \times \pi \times D_{i} \times \Delta T_{lm}} = \frac{3000 \, \text{W}}{500 \, \text{W/m}^2 \text{K} \times \pi \times 0.02 \, \text{m} \times 98.65 \, \text{K}} = 0.968 \, \text{m}
\end{equation}

\subsection*{b. Performance Degradation Due to Fouling}

After three years, fouling has occurred on the water side, reducing the outlet temperature of the water to 65 \( ^{\circ}\)C. Assuming the flow rates and inlet temperatures of both fluids remain unchanged, we now have to calculate the new heat transfer rate \( q_{new} \), the outlet temperature of the oil \( T_{h,o,new} \), the new overall heat transfer coefficient \( U_{new} \), and the fouling factor \( R''_{f,c} \).

First, we update the outlet temperature of the water in the heat transfer rate equation:
\begin{equation}
q_{new} = \dot{m}_{c} C_{p,c} (T_{c,i,new} - T_{c,i}) = 50 \times (65 - 20) = 2250 \, \text{W}
\end{equation}

Next, we solve for the new outlet temperature of the oil:
\begin{equation}
T_{h,o,new} = T_{h,i} - \frac{q_{new}}{\dot{m}_{h} C_{p,h}} = 160 - \frac{2250 \, \text{W}}{150 \, \text{W/K}} = 145 \, ^{\circ}C
\end{equation}

We then recalculate the temperature differences and the LMTD with the new temperatures:
\begin{equation}
\Delta T_{i,new} = T_{h,i} - T_{c,o,new} = 160 - 65 = 95 \, ^{\circ}C
\end{equation}
\begin{equation}
\Delta T_{o,new} = T_{h,o,new} - T_{c,i} = 145 - 20 = 125 \, ^{\circ}C
\end{equation}
\begin{equation}
\Delta T_{lm,new} = \frac{\Delta T_{o,new} - \Delta T_{i,new}}{\ln(\Delta T_{o,new}/\Delta T_{i,new})} = \frac{125 - 95}{\ln(125/95)} = 109.31 \, ^{\circ}C
\end{equation}

Using the new LMTD, we can solve for the new overall heat transfer coefficient:
\begin{equation}
U_{new} = \frac{q_{new}}{A \times \Delta T_{lm,new}} = \frac{2250 \, \text{W}}{\pi \times 0.02 \, \text{m} \times 0.968 \, \text{m} \times 109.31 \, \text{K}} = 338.42 \, \text{W/m}^2\text{K}
\end{equation}

Finally, we can find the water-side fouling factor by rearranging the equation for the overall heat transfer coefficient and solving for \( R''_{f,c} \):
\begin{equation}
R''_{f,c} = \frac{1}{U_{new}} - \frac{1}{U} = \frac{1}{338.42 \, \text{W/m}^2\text{K}} - \frac{1}{500 \, \text{W/m}^2\text{K}} = 0.00095489 \, \text{m}^2\text{K/W}
\end{equation}

This represents the additional thermal resistance due to fouling. It is a measure of the decrease in performance of the heat exchanger due to the buildup of material on the water side that inhibits heat transfer.




\end{document}